\usepackage{xeCJK}
\setCJKmainfont{Songti TC}
\setCJKfamilyfont{kai}{Songti TC}
\newcommand{\kai}{\CJKfamily{kai}} % 用 \kai{使用楷書}

\usepackage{lmodern}
\usepackage{titling}
\usepackage{titlesec}
\usepackage{graphicx}
\usepackage{float}
\usepackage{amsmath, amsthm, amssymb, amsfonts, mathrsfs}
\usepackage{upgreek} % 直立的希臘字母

\usepackage[shortlabels]{enumitem}
\usepackage{setspace}
\renewcommand{\baselinestretch}{1.5}
\usepackage[colorlinks]{hyperref}
\usepackage{fancyvrb}
\usepackage{fancyhdr}
\usepackage{multicol}
\usepackage{caption}
\usepackage{subcaption}
\usepackage{listings}
\usepackage{bm}
\usepackage{bbm} % for \mathbbm{1} (indicator function)
\usepackage{centernot}
\usepackage{mdframed}
\usepackage{xltabular}
\usepackage{tcolorbox}
\tcbuselibrary{skins, breakable}
\usepackage{tikz}
\usetikzlibrary{patterns, arrows.meta, positioning, calc, decorations.pathmorphing, decorations.pathreplacing, decorations.markings, decorations.shapes, angles, quotes, 3d, matrix}
\usepackage[nameinlink, capitalise]{cleveref}
\usepackage{diagbox}
\usepackage{booktabs}
\usepackage{multirow,array}
\usepackage{tabularx}
\usepackage{mathtools}
\usepackage{indentfirst} % 首段 indent
\setlength{\parindent}{2em}
\usepackage{cancel} % 數學式子畫線槓掉
\newcommand\mcancel[2][red]{\setbox0=\hbox{$#2$}%
\rlap{\raisebox{.45\ht0}{\textcolor{#1}{\rule{\wd0}{1pt}}}}#2}
\setcounter{MaxMatrixCols}{20}
\usepackage{datetime2}
\usepackage{paracol}

\makeatletter
\def\tiangan#1{\expandafter\@tiangan\csname c@#1\endcsname}
\def\@tiangan#1{%
  \ifcase#1\or (甲)\or (乙)\or (丙)\or (丁)\or (戊)\or (己)\or (庚)\or (辛)\or (壬)\or (癸)\else(?)\fi
}
\AddEnumerateCounter{\tiangan}{\@tiangan}{(甲)}
\makeatother

% 支援 \wideparen 圓弧記號（若未載入 yhmath，可用此巨集穩定渲染）
\providecommand{\wideparen}[1]{\overset{\frown}{#1}}
\providecommand{\redtext}[1]{{\color{red} #1}}
\providecommand{\bluetext}[1]{{\color{blue} #1}}

\AtBeginDocument{%
 \setlength{\abovedisplayskip}{6pt}
 \setlength{\belowdisplayskip}{6pt}
 \setlength{\abovedisplayshortskip}{4pt}
 \setlength{\belowdisplayshortskip}{4pt}
}

\definecolor{codegreen}{rgb}{0,0.6,0}
\definecolor{codegray}{rgb}{0.5,0.5,0.5}
\definecolor{codepurple}{rgb}{0.58,0,0.82}
\definecolor{backcolour}{rgb}{0.95,0.95,0.92}

\lstdefinestyle{mystyle}{
 backgroundcolor=\color{backcolour},
 commentstyle=\color{codegreen},
 keywordstyle=\color{magenta},
 numberstyle=\tiny\color{codegray},
 stringstyle=\color{codepurple},
 basicstyle=\ttfamily\footnotesize,
 breakatwhitespace=false,
 breaklines=true,
 captionpos=b,
 keepspaces=true,
 numbers=left,
 numbersep=5pt,
 showspaces=false,
 showstringspaces=false,
 showtabs=false,
 tabsize=2
}
\lstset{style=mystyle}

\hypersetup{citecolor=blue, linkcolor=blue, urlcolor=blue}
